\subsection{Characteristic Polynomial Foundations}

\begin{motivation} \label{mot:characteristic-and-minimal-polynomials} \hlt{Two Spectral Polynomials}\\
Eigenvalues alone do not record algebraic repetition, and a matrix can satisfy many polynomial relations.\\
The characteristic polynomial records root multiplicities; the minimal polynomial selects the least-degree monic relation.
\end{motivation}

\begin{definition} \label{def:matrix-characteristic-polynomial} \hlt{Characteristic Polynomial}\\
Let $n\geq1$ and $A\in\operatorname{Mat}(n,n,\mathbb{F})$. The \hlt{characteristic polynomial} of $A$ is $p_A(z)\defeq\det(zI_n-A)\in\mathbb{F}[z]$.
\end{definition}

\begin{theorem} \label{thm:matrix-characteristic-polynomial-basic-coefficients} \hlt{Basic Form of the Characteristic Polynomial}\\
Let $n\geq1$ and $A=[a_{ij}]\in\operatorname{Mat}(n,n,\mathbb{F})$. The following hold:
\begin{enumerate}[label=(\roman*)]
\setlength{\itemsep}{0pt}
\item $p_A(z)=z^n-(\operatorname{tr}A)z^{n-1}+\cdots+(-1)^n\det A$; in particular, $p_A$ is monic of degree $n$.
\item $\operatorname{spec}(A)=\{\lambda\in\mathbb{F}:p_A(\lambda)=0\}$.
\end{enumerate}
\begin{proof}
\begin{enumerate}[label=(\roman*)]
\setlength{\itemsep}{0pt}
\item By Definitions \ref{def:matrix-characteristic-polynomial} and \ref{def:matrix-trace-and-determinant}, $p_A(z)=\sum\limits_{\sigma\in\mathfrak{S}_n}\sign(\sigma)\prod\limits_{j=1}^{n}(z\delta_{\sigma(j),j}-a_{\sigma(j),j})$.\\
A nonidentity permutation moves at least two indices, so its summand has degree at most $n-2$.\\
The identity summand is $\prod\limits_{j=1}^{n}(z-a_{jj})=z^n-(\operatorname{tr}A)z^{n-1}+\cdots$ by Definition \ref{def:matrix-trace-and-determinant}(i).\\
$p_A(0)=\det(-A)=(-1)^n\det A$ by Definition \ref{def:matrix-characteristic-polynomial} and repeated column linearity from Theorem \ref{thm:determinant-basic-laws}(ii). Thus (i) holds.
\item Fix $\lambda\in\mathbb{F}$. By Definition \ref{def:matrix-characteristic-polynomial} and Theorem \ref{thm:determinant-basic-laws}(vii), $p_A(\lambda)=0\Longleftrightarrow\lambda I_n-A$ is singular.\\
By Corollary \ref{cor:matrix-invertible-iff-zero-not-spectrum} and Definition \ref{def:spectrum-and-geometric-multiplicity}(iii), this is equivalent to $\lambda\in\operatorname{spec}(A)$.
\end{enumerate}
\hlt{Structural Commentary}: Permutation degrees determine the coefficients; shifted singularity identifies the root set.\\
\hlt{Key Insight}: The determinant vanishes exactly when the spectral shift has nonzero kernel.
\end{proof}
\end{theorem}

\begin{definition} \label{def:matrix-algebraic-multiplicity} \hlt{Algebraic Multiplicity}\\
Let $n\geq1$, $A\in\operatorname{Mat}(n,n,\mathbb{F})$, and $\lambda\in\mathbb{F}$. The \hlt{algebraic multiplicity} $m_A(\lambda)$ is the multiplicity of $\lambda$ as a root of $p_A$.\\
If $\lambda\notin\operatorname{spec}(A)$, set $m_A(\lambda)\defeq0$. Unqualified multiplicity means algebraic multiplicity.
\end{definition}

Figure \ref{fig:characteristic-polynomial-roots} summarizes the passage from a matrix to its characteristic polynomial and roots.

\begin{illustrationfigure}[H]
  \centering
  \begin{tikzpicture}[
  >=Stealth,
  node distance=0.85cm and 1.0cm,
  box/.style={draw=blue!55!black, fill=blue!5, rounded corners=2pt, align=center, minimum height=0.78cm, minimum width=2.05cm, font=\small},
  root/.style={circle, draw=red!70!black, fill=red!10, inner sep=1.5pt, font=\small},
  arr/.style={->, line width=0.65pt, draw=black!70}
]
  \node[box] (matrix) {$A$};
  \node[box, right=of matrix] (poly) {$p_A$};
  \node[box, right=of poly] (roots) {roots};

  \draw[arr] (matrix) -- (poly);
  \draw[arr] (poly) -- (roots);

  \node[right=0.72cm of roots, inner sep=0pt, font=\small] (openroots) {$\{$};
  \node[root, right=0.10cm of openroots] (l1) {$\lambda_1$};
  \node[root, right=0.24cm of l1] (l2) {$\lambda_2$};
  \node[root, right=0.24cm of l2] (ld) {$\lambda_d$};
  \node[right=0.10cm of ld, inner sep=0pt, font=\small] (closeroots) {$\}$};
  \draw[arr] (roots) -- (openroots);
\end{tikzpicture}

  \caption{The chain $A\mapsto p_A\mapsto\{\lambda_1,\ldots,\lambda_d\}$; $m_A$ retains each root multiplicity.}
  \label{fig:characteristic-polynomial-roots}
\end{illustrationfigure}

\begin{computationmethod} \label{meth:characteristic-equation-eigenspace-computation} \hlt{Characteristic Equation and Eigenspace Computation}\\
Given $n\geq1$ and $A\in\operatorname{Mat}(n,n,\mathbb{F})$, compute its spectral data as follows:
\begin{enumerate}[label=(\roman*)]
\setlength{\itemsep}{0pt}
\item Form $p_A(z)=\det(zI_n-A)$ by Definition \ref{def:matrix-characteristic-polynomial}.
\item Compute $\{\lambda\in\mathbb{F}:p_A(\lambda)=0\}=\operatorname{spec}(A)$ by Theorem \ref{thm:matrix-characteristic-polynomial-basic-coefficients}(ii), recording each $m_A(\lambda)$.
\item $(\forall\,\lambda\in\operatorname{spec}(A))$, row-reduce $A-\lambda I_n$ and compute $E_{\lambda}(A)=\ker(A-\lambda I_n)$.
\item Choose a basis of each $E_{\lambda}(A)$ and return its dimension as the geometric multiplicity.
\end{enumerate}
The output separates algebraic multiplicity from the nullspace dimension.
\end{computationmethod}

\begin{proposition} \label{prop:upper-triangular-characteristic-polynomial} \hlt{Upper-Triangular Eigenvalues}\\
Let $n\geq1$, let $A=[a_{ij}]\in\operatorname{Mat}(n,n,\mathbb{F})$ be upper triangular, and set $J_n\defeq\{1,\ldots,n\}$. Then:
\begin{enumerate}[label=(\roman*)]
\setlength{\itemsep}{0pt}
\item $p_A(z)=\prod\limits_{j=1}^{n}(z-a_{jj})$;
\item $\operatorname{spec}(A)=\{a_{jj}:j\in J_n\}$;
\item $(\forall\,\lambda\in\mathbb{F})$, $m_A(\lambda)=\#\{j\in J_n:a_{jj}=\lambda\}$.
\end{enumerate}
\begin{proof}
\begin{enumerate}[label=(\roman*)]
\setlength{\itemsep}{0pt}
\item Matrix $zI_n-A$ is upper triangular with diagonal entries $z-a_{jj}$. By Theorem \ref{thm:determinant-basic-laws}(v) and Definition \ref{def:matrix-characteristic-polynomial}, $p_A(z)=\prod\limits_{j=1}^{n}(z-a_{jj})$.
\item By (i) and Theorem \ref{thm:matrix-characteristic-polynomial-basic-coefficients}(ii), the roots of $p_A$ and the elements of $\operatorname{spec}(A)$ are exactly the diagonal entries.
\item Fix $\lambda\in\mathbb{F}$. In the factorization from (i), the exponent of $z-\lambda$ is $\#\{j\in J_n:a_{jj}=\lambda\}$.\\
This equals $m_A(\lambda)$ by Definition \ref{def:matrix-algebraic-multiplicity}, including when $\lambda\notin\operatorname{spec}(A)$. Thus (iii) holds.
\end{enumerate}
\hlt{Structural Commentary}: Triangularity reduces the determinant to diagonal factors; repeated factors encode algebraic multiplicity.\\
\hlt{Key Insight}: For an upper-triangular matrix, the entire characteristic root multiset is visible on the diagonal.
\end{proof}
\end{proposition}

\begin{theorem} \label{thm:similarity-preserves-characteristic-polynomial} \hlt{Similarity Preserves Characteristic Data}\\
Let $n\geq1$, let $A,B\in\operatorname{Mat}(n,n,\mathbb{F})$, and let $S\in\operatorname{Mat}(n,n,\mathbb{F})^{\times}$ satisfy $B=SAS^{-1}$. Then:
\begin{enumerate}[label=(\roman*)]
\setlength{\itemsep}{0pt}
\item $p_B=p_A$;
\item $(\forall\,\lambda\in\mathbb{F})$, $m_B(\lambda)=m_A(\lambda)$;
\item $(\forall\,\lambda\in\mathbb{F})$, the restriction $S|_{E_{\lambda}(A)}:E_{\lambda}(A)\to E_{\lambda}(B)$ is an isomorphism, so $\dim E_{\lambda}(B)=\dim E_{\lambda}(A)$.
\end{enumerate}
\begin{proof}
\begin{enumerate}[label=(\roman*)]
\setlength{\itemsep}{0pt}
\item By $zI_n-B=S(zI_n-A)S^{-1}$ and Theorem \ref{thm:determinant-basic-laws}(vi),(viii), $p_B(z)=\det(S)\det(zI_n-A)\det(S^{-1})=p_A(z)$.
\item Equal polynomials have the same roots with the same multiplicities. By (i) and Definition \ref{def:matrix-algebraic-multiplicity}, $(\forall\,\lambda\in\mathbb{F})\,m_B(\lambda)=m_A(\lambda)$.
\item Fix $\lambda\in\mathbb{F}$. If $\mathbf{x}\in E_{\lambda}(A)$, then $(B-\lambda I_n)S\mathbf{x}=S(A-\lambda I_n)\mathbf{x}=\mathbf{0}$, so $S(E_{\lambda}(A))\subseteq E_{\lambda}(B)$. The same calculation for $A=S^{-1}BS$ proves the reverse inclusion. Thus $S|_{E_{\lambda}(A)}$ is an isomorphism onto $E_{\lambda}(B)$, and the dimensions agree.
\end{enumerate}
\hlt{Structural Commentary}: Conjugation preserves the shifted determinant and transports shifted kernels through one coordinate change.\\
\hlt{Key Insight}: Similarity preserves both algebraic multiplicity and geometric multiplicity.
\end{proof}
\end{theorem}

\begin{proposition} \label{prop:block-upper-triangular-characteristic-polynomial} \hlt{Block Upper-Triangular Characteristic Polynomial}\\
Let $r,s\geq1$, let $B\in\operatorname{Mat}(r,r,\mathbb{F})$, $C\in\operatorname{Mat}(r,s,\mathbb{F})$, and $D\in\operatorname{Mat}(s,s,\mathbb{F})$, and set $A\defeq\left(\begin{smallmatrix}B&C\\0&D\end{smallmatrix}\right)$. Then:
\begin{enumerate}[label=(\roman*)]
\setlength{\itemsep}{0pt}
\item $p_A=p_Bp_D$;
\item $(\forall\,\lambda\in\mathbb{F})$, $m_A(\lambda)=m_B(\lambda)+m_D(\lambda)$.
\end{enumerate}
\begin{proof}
\begin{enumerate}[label=(\roman*)]
\setlength{\itemsep}{0pt}
\item The shifted matrix is $zI_{r+s}-A=\left(\begin{smallmatrix}zI_r-B&-C\\0&zI_s-D\end{smallmatrix}\right)$. By Theorem \ref{thm:block-upper-triangular-determinant} and Definition \ref{def:matrix-characteristic-polynomial}, $p_A=p_Bp_D$.
\item Fix $\lambda\in\mathbb{F}$. By (i) and Definition \ref{def:matrix-algebraic-multiplicity}, the multiplicity of $\lambda$ in $p_A$ is the sum of its multiplicities in $p_B$ and $p_D$, including when $\lambda$ is absent from one or both spectra.
\end{enumerate}
\hlt{Structural Commentary}: The determinant splits across the diagonal blocks, so root orders add across the two characteristic factors.\\
\hlt{Key Insight}: The off-diagonal block affects eigenvectors but not the characteristic polynomial.
\end{proof}
\end{proposition}

\begin{theorem} \label{thm:matrix-schur-triangularization-ordered} \hlt{Matrix Schur Triangularisation with Ordered Diagonal}\\
Let $n\geq1$, let $A\in\operatorname{Mat}(n,n,\mathbb{C})$ have eigenvalues $\lambda_1,\ldots,\lambda_n$, counted with algebraic multiplicity and listed in a prescribed order, and let $A\mathbf{x}=\lambda_1\mathbf{x}$ with $\norm{\mathbf{x}}_2=1$. Then:
\begin{enumerate}[label=(\roman*)]
\setlength{\itemsep}{0pt}
\item $(\exists\,U_2\in\operatorname{Mat}(n,n-1,\mathbb{C}))$ such that $[\mathbf{x}\ U_2]$ is unitary and $[\mathbf{x}\ U_2]^*A[\mathbf{x}\ U_2]$ is upper triangular with diagonal $\lambda_1,\ldots,\lambda_n$;
\item if $A\in\operatorname{Mat}(n,n,\mathbb{R})$, $\lambda_1,\ldots,\lambda_n\in\mathbb{R}$, and $\mathbf{x}\in\mathbb{R}^n$, then the matrix in (i) may be chosen real orthogonal.
\end{enumerate}
\begin{proof}
\begin{enumerate}[label=(\roman*)]
\setlength{\itemsep}{0pt}
\item We induct on $n$. The case $n=1$ follows by taking $U=[\mathbf{x}]$.\\
Assume $n>1$. Extend $\mathbf{x}$ to an orthonormal basis by Corollary \ref{cor:orthonormal-list-extends-to-basis}, and write the resulting unitary matrix as $V=[\mathbf{x}\ V_2]$.\\
Since $A\mathbf{x}=\lambda_1\mathbf{x}$, the first column of $V^*AV$ is $\lambda_1\mathbf{e}_1$, so write $V^*AV=\left(\begin{smallmatrix}\lambda_1&\mathbf{b}^*\\0&A'\end{smallmatrix}\right)$ with $\mathbf{b}\in\mathbb{C}^{n-1}$ and $A'\in\operatorname{Mat}(n-1,n-1,\mathbb{C})$. By Theorem \ref{thm:similarity-preserves-characteristic-polynomial}(i) and Proposition \ref{prop:block-upper-triangular-characteristic-polynomial}(i), $p_A(z)=(z-\lambda_1)p_{A'}(z)$. Thus $\lambda_2,\ldots,\lambda_n$ is the prescribed eigenvalue list of $A'$.\\
By the induction hypothesis, choose a unitary $W\in\operatorname{Mat}(n-1,n-1,\mathbb{C})$ such that $T'\defeq W^*A'W$ is upper triangular with diagonal $\lambda_2,\ldots,\lambda_n$. Set $U\defeq V(I_1\oplus W)$. Then $U$ is unitary, its first column is $\mathbf{x}$, and $U^*AU=\left(\begin{smallmatrix}\lambda_1&\mathbf{b}^*W\\0&T'\end{smallmatrix}\right)$ has the required form.
\item We repeat the induction within $\mathbb{R}^n$. The case $n=1$ is immediate. For $n>1$, extend the real unit vector $\mathbf{x}$ to a real orthonormal basis by Corollary \ref{cor:orthonormal-list-extends-to-basis}. The resulting real orthogonal $V=[\mathbf{x}\ V_2]$ satisfies $V^{\mathsf{T}}AV=\left(\begin{smallmatrix}\lambda_1&\mathbf{b}^{\mathsf{T}}\\0&A'\end{smallmatrix}\right)$ with real blocks.\\
By the factorization in (i), the real matrix $A'$ has the prescribed real eigenvalue list $\lambda_2,\ldots,\lambda_n$.\\
By the real induction hypothesis, choose a real orthogonal $W$ with $W^{\mathsf{T}}A'W$ upper triangular in that order. Thus $U\defeq V(I_1\oplus W)$ is real orthogonal, begins with $\mathbf{x}$, and has the required triangularization.
\end{enumerate}
\hlt{Structural Commentary}: Schur induction isolates the prescribed eigenline and passes the remaining ordered roots to the trailing block.\\
\hlt{Key Insight}: The first eigenvector fixes the first diagonal entry and leaves the remaining multiplicities for the next step.
\end{proof}
\end{theorem}

Figure \ref{fig:schur-ordered-triangularization} summarizes the induction by isolating the prescribed first eigenline.

\begin{illustrationfigure}[H]
  \centering
  \begin{tikzpicture}[
  >=Stealth,
  every node/.style={font=\small},
  block/.style={draw=gray!55, fill=gray!5, rounded corners=2pt, inner sep=6pt, align=center, minimum height=1.0cm},
  focus/.style={draw=blue!65!black, fill=blue!7, rounded corners=2pt, inner sep=6pt, align=center, minimum height=1.0cm},
  arr/.style={->, line width=0.75pt, blue!70!black}
]
  \node[focus] (line) at (0,0)
    {$A\mathbf{x}=\lambda_1\mathbf{x}$};
  \node[block] (block) at (3.3,0)
    {$V^*AV=\begin{bmatrix}\lambda_1&*\\0&A'\end{bmatrix}$};
  \node[block] (tri) at (6.9,0)
    {$U^*AU=T$};

  \draw[arr] (line) -- (block);
  \draw[arr] (block) -- (tri);
\end{tikzpicture}

  \caption{Schur induction: isolate an eigenline, then triangularise the compressed block.}
  \label{fig:schur-ordered-triangularization}
\end{illustrationfigure}

\begin{corollary} \label{cor:eigenvalues-of-inverse-matrix} \hlt{Eigenvalues of the Inverse}\\
Let $n\geq1$ and let $A\in\operatorname{Mat}(n,n,\mathbb{C})$ be invertible with eigenvalues $\lambda_1,\ldots,\lambda_n$, counted with algebraic multiplicity.\\
Then $A^{-1}$ has eigenvalues $\lambda_1^{-1},\ldots,\lambda_n^{-1}$, counted with algebraic multiplicity.
\begin{proof}
By Corollary \ref{cor:matrix-invertible-iff-zero-not-spectrum}, $(\forall\,j\in\{1,\ldots,n\})\,\lambda_j\neq0$. Choose $\mathbf{0}\neq\mathbf{y}\in E_{\lambda_1}(A)$ and set $\mathbf{x}\defeq\frac{\mathbf{y}}{\norm{\mathbf{y}}_2}$, so $\norm{\mathbf{x}}_2=1$.\\
By Theorem \ref{thm:matrix-schur-triangularization-ordered}(i), $A=UTU^*$ for a unitary $U$ and an upper-triangular $T$ with diagonal $\lambda_1,\ldots,\lambda_n$.\\
Hence $A^{-1}=UT^{-1}U^*$, while $T^{-1}$ is upper triangular with diagonal $\lambda_1^{-1},\ldots,\lambda_n^{-1}$ by Theorem \ref{thm:inverse-upper-triangular}.\\
By Proposition \ref{prop:upper-triangular-characteristic-polynomial}(iii) and Theorem \ref{thm:similarity-preserves-characteristic-polynomial}(ii), this reciprocal diagonal list is the eigenvalue multiset of $A^{-1}$.\\
\hlt{Structural Commentary}: Schur form turns inversion into inversion of a triangular matrix, whose diagonal entries invert individually.\\
\hlt{Key Insight}: Algebraic multiplicities are preserved when the nonzero spectral multiset is sent through $\lambda\mapsto\lambda^{-1}$.
\end{proof}
\end{corollary}

\begin{theorem} \label{thm:geometric-multiplicity-bounded-by-algebraic-multiplicity} \hlt{Geometric Multiplicity Is at Most Algebraic Multiplicity}\\
Let $n\geq1$, let $A\in\operatorname{Mat}(n,n,\mathbb{F})$, and let $\lambda\in\operatorname{spec}(A)$. Then $\dim E_{\lambda}(A)\leq m_A(\lambda)$.
\begin{proof}
Set $k\defeq\dim E_{\lambda}(A)$ and choose a basis $(\mathbf{x}_1,\ldots,\mathbf{x}_k)$ of $E_{\lambda}(A)$.\\
If $k=n$, then $E_{\lambda}(A)=\mathbb{F}^n$, so $A=\lambda I_n$. By Proposition \ref{prop:upper-triangular-characteristic-polynomial}(i),(iii), $p_A(z)=(z-\lambda)^n$ and $m_A(\lambda)=n=k$. Suppose $k<n$.\\
Extend eigenspace basis to basis of $\mathbb{F}^n$ by Theorem \ref{thm:hamel-basis-extension}, let $S$ be invertible matrix whose columns are resulting basis vectors.\\
Since the first $k$ columns are $\lambda$-eigenvectors, write $S^{-1}AS=\left(\begin{smallmatrix}\lambda I_k&C\\0&D\end{smallmatrix}\right)$ with $C\in\operatorname{Mat}(k,n-k,\mathbb{F})$ and $D\in\operatorname{Mat}(n-k,n-k,\mathbb{F})$.\\
By Theorem \ref{thm:similarity-preserves-characteristic-polynomial}(i) and Proposition \ref{prop:block-upper-triangular-characteristic-polynomial}(i), $p_A(z)=(z-\lambda)^kp_D(z)$. Hence $m_A(\lambda)\geq k=\dim E_{\lambda}(A)$ by Definition \ref{def:matrix-algebraic-multiplicity}.\\
\hlt{Structural Commentary}: An eigenspace-first basis exposes $(z-\lambda)^k$, including the endpoint $k=n$.\\
\hlt{Key Insight}: Geometric multiplicity supplies a visible lower bound on the order of the corresponding characteristic root.
\end{proof}
\end{theorem}

\begin{theorem} \label{thm:diagonalizable-iff-algebraic-geometric-multiplicities-match} \hlt{Diagonalizability and Multiplicity}\\
Let $n\geq1$ and let $A\in\operatorname{Mat}(n,n,\mathbb{C})$. Then $A$ is similar to a diagonal matrix if and only if $(\forall\,\lambda\in\operatorname{spec}(A))\,\dim E_{\lambda}(A)=m_A(\lambda)$.
\begin{proof}
\begin{enumerate}[label=\roman*.]
\setlength{\itemsep}{0pt}
\item[$(\Rightarrow)$] Suppose $A$ is similar to a diagonal matrix $D$. By Theorem \ref{thm:similarity-preserves-characteristic-polynomial}(ii),(iii), $(\forall\,\lambda\in\mathbb{C})\,m_A(\lambda)=m_D(\lambda)$ and $\dim E_{\lambda}(A)=\dim E_{\lambda}(D)$.\\
By Proposition \ref{prop:upper-triangular-characteristic-polynomial}(iii), $m_D(\lambda)$ counts the diagonal entries equal to $\lambda$, and the corresponding standard basis vectors form a basis of $E_{\lambda}(D)$. Thus the required equality holds.
\item[$(\Leftarrow)$] Assume the stated equality and let $T_A:\mathbb{C}^n\to\mathbb{C}^n$ be defined by $T_A(\mathbf{v})=A\mathbf{v}$. By Theorem \ref{thm:complex-upper-triangularization}, $T_A$ has an upper-triangular matrix $B$ relative to a basis. By Corollary \ref{cor:similar-iff-same-linear-operator} and Theorem \ref{thm:similarity-preserves-characteristic-polynomial}(ii), $A$ and $B$ have the same algebraic multiplicities.\\
List the distinct eigenvalues as $\mu_1,\ldots,\mu_d$. By Proposition \ref{prop:upper-triangular-characteristic-polynomial}(ii),(iii), the $n$ diagonal positions of $B$ are partitioned by these eigenvalues and $\sum\limits_{i=1}^{d}m_A(\mu_i)=n$. Thus the hypothesis implies $\sum\limits_{i=1}^{d}\dim E_{\mu_i}(T_A)=\sum\limits_{i=1}^{d}\dim E_{\mu_i}(A)=n$.\\
By Theorem \ref{thm:conditions-equivalent-to-diagonalizability}(v),(i), $T_A$ is diagonalizable. Its matrices in the standard basis and in a diagonalizing basis are therefore similar by Corollary \ref{cor:similar-iff-same-linear-operator}, so $A$ is similar to a diagonal matrix.
\end{enumerate}
\hlt{Structural Commentary}: Triangularization accounts for the complete algebraic-multiplicity total, while the hypothesis converts it into a full eigenspace dimension count.\\
\hlt{Key Insight}: A complex matrix is diagonalizable exactly when no algebraic multiplicity is lost as an eigenspace dimension.
\end{proof}
\end{theorem}

\begin{proposition} \label{prop:permutation-similarity-reorders-diagonal} \hlt{Permutation Similarity Reorders Diagonal Entries}\\
Let $n\geq1$, let $\lambda_1,\ldots,\lambda_n\in\mathbb{C}$, and let $\sigma\in\mathfrak{S}_n$. Set $\Lambda\defeq\operatorname{diag}(\lambda_1,\ldots,\lambda_n)$ and $P\defeq[\mathbf{e}_{\sigma(1)}\ \cdots\ \mathbf{e}_{\sigma(n)}]$. Then $P^*\Lambda P=\operatorname{diag}(\lambda_{\sigma(1)},\ldots,\lambda_{\sigma(n)})$.
\begin{proof}
For $i,j\in\{1,\ldots,n\}$, the $(i,j)$ entry of $P^*\Lambda P$ is $\mathbf{e}_{\sigma(i)}^*\Lambda\mathbf{e}_{\sigma(j)}=\lambda_{\sigma(j)}\mathbf{e}_{\sigma(i)}^*\mathbf{e}_{\sigma(j)}=\lambda_{\sigma(j)}\delta_{ij}$.\\
Hence all off-diagonal entries vanish and the $j$th diagonal entry is $\lambda_{\sigma(j)}$, which proves the matrix identity.\\
\hlt{Structural Commentary}: The entry calculation tracks both the row and column reordering induced by the same permutation.\\
\hlt{Key Insight}: Permutation similarity reorders a diagonal list without changing its entries.
\end{proof}
\end{proposition}

\begin{corollary} \label{cor:normal-matrices-same-spectrum-unitarily-similar} \hlt{Normal Matrices with the Same Spectrum Are Unitarily Similar}\\
Let $n\geq1$. If $A,B\in\operatorname{Mat}(n,n,\mathbb{C})$ are normal and have the same eigenvalue multiset, then $A$ and $B$ are unitarily similar.
\begin{proof}
By Theorem \ref{thm:complex-spectral-theorem}(i),(iii), write $A=U\Lambda U^*$ and $B=V\Gamma V^*$ with $U,V$ unitary and $\Lambda,\Gamma$ diagonal.\\
Since the eigenvalue multisets agree, choose $\sigma\in\mathfrak{S}_n$ that matches the diagonal of $\Lambda$ to that of $\Gamma$, and set $P\defeq[\mathbf{e}_{\sigma(1)}\ \cdots\ \mathbf{e}_{\sigma(n)}]$.\\
By Proposition \ref{prop:permutation-similarity-reorders-diagonal}, $\Gamma=P^*\Lambda P$. Set $Q\defeq UPV^*$.\\
Then $Q$ is unitary by Proposition \ref{prop:algebra-of-unitary-matrices}, and $Q^*AQ=VP^*U^*(U\Lambda U^*)UPV^*=V\Gamma V^*=B$. Thus $A$ and $B$ are unitarily similar.\\
\hlt{Structural Commentary}: Spectral diagonalizations reduce the comparison to a unitary permutation of equal eigenvalue multisets.\\
\hlt{Key Insight}: For normal matrices, the eigenvalue multiset completely determines the unitary-similarity class.
\end{proof}
\end{corollary}

\begin{definition} \label{def:operator-characteristic-polynomial} \hlt{Characteristic Polynomial of an Operator}\\
Let $\mathcal{V}\neq\{\mathbf{0}\}$ be finite-dimensional over $\mathbb{F}$, let $n\defeq\dim(\mathcal{V})$, and let $T\in\LL(\mathcal{V})$.\\
For a basis $\beta$ of $\mathcal{V}$, set $q_{T,\beta}(z)\defeq p_{[T]_{\beta\leftarrow\beta}}(z)=\det(zI_n-[T]_{\beta\leftarrow\beta})$.\\
This value is independent of $\beta$ by Theorems \ref{thm:change-of-basis-formula} and \ref{thm:similarity-preserves-characteristic-polynomial}(i). It is the \hlt{characteristic polynomial} of $T$ and is denoted by $q_T$.
\end{definition}

\begin{proposition} \label{prop:operator-characteristic-polynomial-degree-roots} \hlt{Degree and Roots of $q_T$}\\
Let $\mathcal{V}\neq\{\mathbf{0}\}$ be finite-dimensional over $\mathbb{F}$, let $T\in\LL(\mathcal{V})$, and let $q_T$ be the characteristic polynomial of $T$.
\begin{enumerate}[label=(\roman*)]
\setlength{\itemsep}{0pt}
\item $\deg q_T=\dim(\mathcal{V})$.
\item The roots of $q_T$ are precisely the eigenvalues of $T$.
\end{enumerate}
\begin{proof}
Choose a basis $\beta$ of $\mathcal{V}$ and set $A=[T]_{\beta\leftarrow\beta}$.
\begin{enumerate}[label=(\roman*)]
\setlength{\itemsep}{0pt}
\item By Definition \ref{def:operator-characteristic-polynomial}, $q_T=p_A$. Hence $\deg q_T=\dim(\mathcal{V})$ by Theorem \ref{thm:matrix-characteristic-polynomial-basic-coefficients}(i).
\item Fix $\lambda\in\mathbb{F}$. By Proposition \ref{prop:matrix-representation-action}, $(\forall\,\mathbf{v}\in\mathcal{V})\,[T\mathbf{v}]_{\beta}=A[\mathbf{v}]_{\beta}$.\\
By Proposition \ref{prop:coordinate-map-isomorphism}(ii),(iii), $\lambda\in\operatorname{spec}(T)\Longleftrightarrow\lambda\in\operatorname{spec}(A)$. By Theorem \ref{thm:matrix-characteristic-polynomial-basic-coefficients}(ii), the latter set is the root set of $p_A=q_T$.
\end{enumerate}
\hlt{Structural Commentary}: Coordinates transfer polynomial degree and the nonzero eigenvector equation between $T$ and its matrix.\\
\hlt{Key Insight}: The operator characteristic polynomial has the expected dimension and detects exactly the operator eigenvalues.
\end{proof}
\end{proposition}
